% SedSAT3 – Response to the USGS code review (S. Westenbroek, October 2026)
% Build: pdflatex SedSAT3_code_review_response.tex  (run twice)
\documentclass[11pt]{article}
\usepackage[letterpaper,margin=1in]{geometry}
\usepackage[T1]{fontenc}
\usepackage{lmodern}
\usepackage{microtype}
\usepackage[utf8]{inputenc}
\usepackage{parskip}
\usepackage{xcolor}
\usepackage{enumitem}
\usepackage{booktabs}
\usepackage{longtable}
\usepackage{tabularx}
\usepackage[most]{tcolorbox}
\usepackage{fancyhdr}
\usepackage[hidelinks]{hyperref}

\definecolor{usgsgreen}{RGB}{0,100,60}
\definecolor{commentbg}{RGB}{242,245,248}
\definecolor{todo}{RGB}{190,30,30}
\definecolor{mine}{RGB}{20,70,170}

% ---------------------------------------------------------------- macros
\newcommand{\code}[1]{\texttt{#1}}
\newcommand{\commit}[1]{\texttt{#1}}
\newcommand{\todo}[1]{\textcolor{todo}{\textbf{[#1]}}}   % for Cara / Allen / the release coordinator
\newcommand{\mine}[1]{\textcolor{mine}{\textbf{[Arash: #1]}}}   % programming items for the author
\newcounter{cmt}
\newcounter{grp}
\renewcommand{\thecmt}{\thegrp.\arabic{cmt}}
\newcommand{\group}[1]{\stepcounter{grp}\setcounter{cmt}{0}\section*{\thegrp.\quad #1}\addcontentsline{toc}{section}{\thegrp.\quad #1}}
% \reviewer{source}{comment text}
\newenvironment{reviewer}[1]{%
  \refstepcounter{cmt}%
  \begin{tcolorbox}[enhanced,breakable,colback=commentbg,colframe=usgsgreen!60,
    boxrule=0.4pt,left=6pt,right=6pt,top=4pt,bottom=4pt,
    title={\small\textbf{Comment \thecmt}\hfill\textit{#1}},
    coltitle=black,colbacktitle=usgsgreen!15,fonttitle=\small]}%
  {\end{tcolorbox}}
\newenvironment{response}{\par\noindent\textbf{Response.}\ }{\par\medskip}
\newcommand{\status}[1]{\par\noindent\textit{Status:} #1\par}

\pagestyle{fancy}
\fancyhf{}
\lhead{\small SedSAT3 — response to code review}
\rhead{\small \thepage}
\renewcommand{\headrulewidth}{0.3pt}

\begin{document}
\sloppy
\setlength{\emergencystretch}{3em}

% ---------------------------------------------------------------- letter
\begin{tcolorbox}[colback=todo!4,colframe=todo!70,boxrule=0.6pt]
\textbf{Draft for Cara Peterman-Phipps and Allen Gellis.} Items in
\textcolor{todo}{\textbf{red}} need your input or the release coordinator's;
items in \textcolor{mine}{\textbf{blue}} are programming tasks Arash is
completing. All marked items and this note will be removed before the letter
goes to Steve.
\end{tcolorbox}

\begin{flushleft}
Arash Massoudieh\\
Department of Civil and Environmental Engineering\\
The Catholic University of America\\
\href{mailto:massoudieh@cua.edu}{massoudieh@cua.edu}

\medskip
\today

\medskip
Stephen M.\ Westenbroek\\
U.S. Geological Survey\\
via Cara Peterman-Phipps, U.S. Geological Survey
\end{flushleft}

\textbf{Re: Code review of SedSAT3 (Sediment Source Assessment Tool, version 3), IPDS IP-190994}

Dear Steve,

Thank you for the careful and generous review of SedSAT3, and for the
additional material you sent on October 16. Both were very useful. Your
notes on how to read the Kiro-generated documents (as a triage guide, with
the line numbers worth checking and the prose taken with a grain of salt)
were exactly the right framing, and I have tried to follow it: every
concrete finding was checked against the code, real defects were fixed, and
findings that turned out not to apply are explained rather than worked
around.

All changes are in merge request
\href{https://code.usgs.gov/cphipps/SedimentSourceAssessmentTool3/-/merge_requests/4}{!4}
on code.usgs.gov (branch \code{review-fixes-oct2026}), organised as one commit
per area so they can be reviewed separately:

\begin{center}
\small
\begin{tabularx}{\linewidth}{@{}lX@{}}
\toprule
Commit & Content\\
\midrule
\commit{ac1dcd8} & Review notes, test projects, and a regression baseline recorded before any change\\
\commit{41393e0} & Licensing, attribution and USGS metadata\\
\commit{ac039a2} & The confirmed code defects\\
\commit{cdfa4ea} & Seeded, thread-independent random numbers\\
\commit{68be7ca} & Screening of every imported value\\
\commit{2f8c463} & Unit tests, regression tests, GitLab CI, optional GA, build cleanup, warnings\\
\commit{538ed73} & Response document (\code{docs/\allowbreak{}review/\allowbreak{}REVIEW\_\allowbreak{}RESPONSE.\allowbreak{}md})\\
\commit{follow-up} & Follow-up to your October 16 documents: further fixes and tests (Sections 7--11)\\
\bottomrule
\end{tabularx}
\end{center}

In brief:
\begin{itemize}[nosep]
\item All non-USGS license and attribution notices are gone; the repository
  no longer depends on or points to any personal repository; LICENSE,
  DISCLAIMER, code.json, CITATION.cff and a third-party license inventory are
  in place. The license itself, the DOI and the author list await the release
  coordinator (Section~1).
\item Every code defect in your review was confirmed and fixed, plus about a
  dozen further defects found while doing so (Sections~2 and~6).
\item The MCMC now takes a user-set random seed, records it, and gives
  identical results for the same seed regardless of the number of threads
  (Section~3).
\item Every imported value is screened by a single parser; NaN and Inf are
  refused at import with the cell location (Section~4).
\item There are now 49 unit tests (including the NaN / negative / illegal
  value tests you suggested), a regression suite on two example projects,
  and a GitLab CI pipeline running the tests, clang-tidy and cppcheck
  (Section~5).
\end{itemize}

Below, each of your comments is reproduced (quoted where short, summarised
where long) and followed by my response. Comments are grouped by topic; the
source of each is given at the right of its heading: \emph{Email 10/15} and
\emph{Email 10/16} are your two emails to Cara, \emph{Summary} is
\code{CODE\_\allowbreak{}REVIEW\_\allowbreak{}SUMMARY\_\allowbreak{}15OCT2026.\allowbreak{}md}, \emph{Tests} is
\code{TEST\_\allowbreak{}RESULTS.\allowbreak{}md}, \emph{Observations} is
\code{CODE\_\allowbreak{}REVIEW\_\allowbreak{}OBSERVATIONS.\allowbreak{}md}, and \emph{Linter} is
\code{linter\_\allowbreak{}output.\allowbreak{}txt}. File and line references are to the code after the
fixes (commit \commit{538ed73}).

Thank you again for your time; I would be very glad to have you take another
look once the licensing and ownership questions are settled.

Sincerely,

\medskip
Arash Massoudieh

\newpage
\tableofcontents
\newpage

% ====================================================================
\group{Ownership, licensing, attribution and citation}

\begin{reviewer}{Email 10/15; Summary, threshold question}
In its current form too much of the code and documentation references
non-USGS resources. ``I'd encourage you to clarify the issue of ownership
and commit to either pushing full on for USGS support (i.e.\ documentation and
code and issues migrated to a USGS supported platform) or just continue to
develop as is without trying to rebrand as USGS software.''
\end{reviewer}
\begin{response}
\todo{Confirm wording with Cara and the release coordinator.} We intend
SedSAT3 to be a USGS software release, as the successor of the USGS Sed\_SAT
tool (OFR 2017-1062). The repository on code.usgs.gov is now the
authoritative copy: code, documentation and issue tracking live there, the
README, code.json and CITATION.cff point only to it (placeholder until the
final group path is assigned), and no file refers to a personal repository.
Following your suggestion, any further development on my side will be
submitted as merge requests against the USGS repository.
\end{response}
\status{Repository side done (\commit{41393e0}); institutional decision with the release coordinator.}

\begin{reviewer}{Email 10/15; Summary, blocker 1}
Several files carry OpenHydroQual / GNU AGPL-3.0 headers with a ``purchase a
commercial license, contact \ldots enviroinformatics.co'' clause, contradicting
the Utilities submodule's own Apache-2.0 license and far from the USGS
CC0 preference.
\end{reviewer}
\begin{response}
These files (\code{TimeSeries.\allowbreak{}h/\allowbreak{}.\allowbreak{}hpp}, \code{TimeSeriesSet.\allowbreak{}h/\allowbreak{}.\allowbreak{}hpp}) were
written by me for another project and reused here. The \code{Utilities}
submodule has been removed and the 23 files SedSAT3 actually builds were
brought into the repository (\code{Utilities/\allowbreak{}README.\allowbreak{}md} records where they
came from). The AGPL / commercial headers were replaced by the same two-line
USGS header that every first-party source file now carries:
\begin{quote}\small\ttfamily
// SedSAT3 – Sediment Source Assessment Tool, version 3\\
// U.S. Geological Survey. See LICENSE.md and DISCLAIMER.md in the repository root.
\end{quote}
As the author I release these copies for SedSAT3 under the repository
license; I am happy to provide a signed statement to that effect if the
release coordinator needs one. Third-party code (QXlsx, QCustomPlot) was not
touched.
\end{response}
\status{Fixed (\commit{41393e0}).}

\begin{reviewer}{Email 10/15; Summary, blocker 2}
The repository is GPL-3.0 because it links GSL, QCustomPlot and Qt Charts,
outside the USGS preference (CC0/MIT/Unlicense). Rather than re-engineering,
consider the SWB2 approach (a third-party license inventory); the USGS
package pfdf was deliberately released under GPL.
\end{reviewer}
\begin{response}
I followed your SWB2 suggestion rather than replacing libraries.
\code{THIRD\_\allowbreak{}PARTY\_\allowbreak{}LICENSES.\allowbreak{}md} lists every dependency (Qt modules, Qt
Charts, QCustomPlot, QXlsx, GSL, Armadillo, LAPACK/BLAS, the OpenMP and MSVC
runtimes), how SedSAT3 uses each, its license, and where it comes from, and
notes that a binary package linking the GPL components is distributed under
GPL-compatible terms. \code{LICENSE.\allowbreak{}md} currently holds a placeholder with
the USGS CC0 text ready. \todo{License decision: CC0 1.0, MIT, or — following
the pfdf precedent — GPL-3.0, which would remove the dependency question
entirely. To be decided with Cara and the release coordinator; I plan to ask
Jonathan King what the pfdf GPL release required.}
\end{response}
\status{Inventory done (\commit{41393e0}); license choice open.}

\begin{reviewer}{Summary, blocker 3}
The canonical identity, the \code{Utilities} submodule and the QXlsx fork
live under the author's personal GitHub, which points toward non-USGS
primary authorship. Raise with the release coordinator.
\end{reviewer}
\begin{response}
The \code{Utilities} submodule is gone (see Comment 1.2) and the QXlsx
submodule now points at its public upstream (\code{QtExcel/\allowbreak{}QXlsx}), not a
personal fork. Authorship is recorded in \code{CITATION.\allowbreak{}cff} and the README
citation \todo{author list and order to be confirmed: Massoudieh, A.;
Peterman-Phipps, C.; Gellis, A.C.; \ldots}.
\end{response}
\status{Repository side done (\commit{41393e0}, \commit{2f8c463}); authorship for the release coordinator.}

\begin{reviewer}{Email 10/15}
The README points to GitHub for the code and the user manual, and the
suggested citation (``Massoudieh, A., et al.\ (2023) \ldots
https://github.com/\ldots'') ``fails on many levels (a USGS product needs to
cite a DOI to a durable, USGS-approved digital repository; all contributors
need to be listed).''
\end{reviewer}
\begin{response}
The README was rewritten: no GitHub links; the repository is the
code.usgs.gov project; the manual points to its USGS location
\todo{manual URL/DOI}; the citation is in USGS software-release form with
the full author list and the release DOI \todo{10.5066/XXXXXXX}. The same
metadata is in \code{CITATION.\allowbreak{}cff} and \code{code.\allowbreak{}json}. The README also
gained sections on the disclaimer, license, build instructions, build
options, tests, and running the portable Windows build.
\end{response}
\status{Fixed with placeholders (\commit{41393e0}).}

\begin{reviewer}{Summary, should address}
Metadata (\code{code.\allowbreak{}json}, README, \code{createdeb.\allowbreak{}sh}, submodules) points at
personal GitHub; the README's Windows installer link (sedsat.org) serves a
different-lineage artifact (``CMBSource'' by ``EnviroInformatics'').
\end{reviewer}
\begin{response}
\code{code.\allowbreak{}json} was rewritten for code.usgs.gov (and no longer lists the
genetic algorithm as a method); \code{createdeb.\allowbreak{}sh} now only builds and
checksums the Debian package, with no GitHub publishing; the README no
longer links to sedsat.org. In the installer project the manufacturer is now
``U.S. Geological Survey'' and its URLs point to the USGS repository. The
Windows package reviewers receive will be the portable build made from this
repository (Section~7).
\end{response}
\status{Fixed (\commit{41393e0}).}

\begin{reviewer}{Summary, optional / polish}
No per-file license headers; no shipped dependency-license (NOTICE)
inventory. The README is truncated and reads as adapted from the predecessor
``CMBSource'' project.
\end{reviewer}
\begin{response}
Every first-party source file now carries the USGS header, and
\code{THIRD\_\allowbreak{}PARTY\_\allowbreak{}LICENSES.\allowbreak{}md} is the dependency inventory. The README was
rewritten from scratch (Comment 1.5).
\end{response}
\status{Fixed (\commit{41393e0}).}

% ====================================================================
\group{Code defects}

\begin{reviewer}{Email 10/15; Summary; Tests, Test 7}
Distribution Fitting terminates the application with an unhandled
\code{std::\allowbreak{}out\_\allowbreak{}of\_\allowbreak{}range} (debugger-traced to \code{TimeSeriesSet::\allowbreak{}operator[]}
from \code{ExecuteDistributionFitting}); suspected cause: a non-existent
\code{TimeSeriesSet::\allowbreak{}append(series, name)} overload, so fit results are stored
unnamed.
\end{reviewer}
\begin{response}
Confirmed and fixed, though the root cause differs from Kiro's diagnosis: the
\code{append(series, name)} overload exists in this version. The fault was
in \code{ConcentrationSet::\allowbreak{}CreateFittedDistribution}, which assigned a number
to the ``Observed'' series (\code{result["Observed"] = maxC/\allowbreak{}2}). Because
\code{TimeSeries(size\_\allowbreak{}t)} was not \code{explicit}, the number was silently
converted into an empty, unnamed series, and the later lookup of
``Observed'' threw. Three further faults in the same path were fixed: the
log-normal series was silently dropped by a duplicate-name check and
\code{setname(2,...)} then indexed past the end; empty or constant data made
the fitting loop infinite (zero step); and the Kolmogorov--Smirnov statistic
indexed a series that might not exist. \code{TimeSeries(size\_\allowbreak{}t)} is now
\code{explicit}, so this class of silent conversion cannot recur; the analysis
checks the group, constituent, sample count and positivity and reports a
clear message; and \code{MainWindow::\allowbreak{}Execute} now catches any exception from
any analysis and reports it in a dialog instead of terminating.
\end{response}
\status{Fixed (\commit{ac039a2}). Tests: \code{distributionFitting\_\allowbreak{}degenerateSets},
\code{distributionFitting\_\allowbreak{}conductor}; regression (all distribution-fitting steps now succeed).}

\begin{reviewer}{Email 10/15; Summary; Tests, G1 and G2}
The ``Include all elements'' checkbox has inverted logic (checking it
excludes all tracers), and its label flips between ``Include'' and
``Exclude'', which is what produced the inversion. ``Just ensure that the
behavior of the checkboxes matches your intentions.''
\end{reviewer}
\begin{response}
Fixed as you suggest: the label is now fixed (``Include all elements''),
checked means every element is included, and the box follows the
per-element settings (it clears when any element is excluded in the table).
\code{IncludeExcludeAllElements} has no other caller. While writing the test
I also found that the element table read the text ``No'' as \emph{include}
(\code{QVariant("No").\allowbreak{}toBool()} is true); fixed.
\end{response}
\status{Fixed (\commit{ac039a2}, \commit{2f8c463}). Tests: \code{includeAll\_\allowbreak{}dataLevel}, \code{includeAll\_\allowbreak{}checkBox}.}

\begin{reviewer}{Tests, G3}
Toggling a single element appears to require several clicks.
\end{reviewer}
\begin{response}
Fixed: one click on the ``Include in analysis'' cell now toggles it. The
cause of the extra clicks was a bug: the cell opened a check-box editor that
started ticked even for excluded elements (it read the text ``No'' as true),
so the first click often changed nothing.
\end{response}
\status{Fixed. Test: \code{includeColumn\_\allowbreak{}singleClickToggles}.}

\begin{reviewer}{Email 10/15; Summary; Tests, Test 4c}
Missing \code{return} on the clean path of
\code{Conductor::\allowbreak{}CheckNegativeElements} (both overloads) — undefined behaviour,
currently masked on MSVC.
\end{reviewer}
\begin{response}
This had already been fixed on \code{main} (\commit{238dd69}) after the
reviewed snapshot. I then searched for the same defect everywhere: GCC
\code{-Wreturn-type} found six more functions that could fall off their end,
all now fixed. To make this impossible in future, \code{-Werror=return-type}
(GCC/Clang) and \code{/\allowbreak{}we4715} (MSVC) are set in every build file.
While testing this guard I found a more serious defect in it:
\code{NegativeValueCheck} (and \code{GrandMean} and \code{LumpAllProfileSets})
stopped before the last group in alphabetical order, assuming it was the
target. With target ``Mixture'', the source ``Pasture'' was never checked and
the target was checked instead. All three now skip the target by name.
\end{response}
\status{Fixed (\commit{ac039a2}, \commit{2f8c463}). Tests: \code{negativeCheck\_\allowbreak{}clean},
\code{negativeCheck\_\allowbreak{}negative}, \code{negativeCheck\_\allowbreak{}sourceAfterTarget}.}

\begin{reviewer}{Tests, preprocessing note and Test 4 addendum}
The negative-value guard rejects legitimately negative isotopic tracers
(e.g.\ $\delta^{13}$C) with no transform path, and its message is
misleading when the negativity is produced by the organic-matter / size
correction rather than present in the input.
\end{reviewer}
\begin{response}
The message now distinguishes the two cases: when a value is negative only
after the correction, it is reported as ``(after the organic matter and size
correction; the imported values are positive)''. Isotopic tracers such as
$\delta^{13}$C are meant to be given the \emph{Isotope} role (with a base
element and standard ratio), under which they are handled through their
ratio rather than log-transformed; the README now says so in the quick
start. Whether an OM/size correction that produces negative concentrations
is acceptable is a scientific question I have flagged for the domain review.
\end{response}
\status{Improved (\commit{538ed73}); domain question handed off.}

\begin{reviewer}{Email 10/15; Summary}
Null dereference in \code{CBinary::\allowbreak{}operator[]} (\code{src/\allowbreak{}GA/\allowbreak{}Binary.\allowbreak{}cpp}).
\end{reviewer}
\begin{response}
Fixed: an out-of-range index now throws \code{std::\allowbreak{}out\_\allowbreak{}of\_\allowbreak{}range} with the
index and size. The same deliberate null dereference was the out-of-range
path of \code{CVector::\allowbreak{}operator[]} and \code{CVector\_\allowbreak{}arma::\allowbreak{}operator[]}, which
are used throughout the analyses; those now throw too, so an indexing
mistake anywhere produces a message (through the GUI's exception handler)
rather than a silent crash.
\end{response}
\status{Fixed (\commit{ac039a2}, \commit{2f8c463}). Test: \code{binary\_\allowbreak{}indexing}.}

\begin{reviewer}{Email 10/15; Summary}
Writes to a closed \code{FILE*} in \code{src/\allowbreak{}GA/\allowbreak{}GA.\allowbreak{}cpp} (lines 308--309).
\end{reviewer}
\begin{response}
\code{src/\allowbreak{}GA/\allowbreak{}GA.\allowbreak{}cpp} turned out to be a stale copy of the GA that no build
compiled (it includes a \code{StringOP.\allowbreak{}h} that exists nowhere — the
``missing header'' you noted); it was removed. The GA code that is compiled,
\code{include/\allowbreak{}GA/\allowbreak{}GA.\allowbreak{}hpp}, now opens its files through a \code{unique\_\allowbreak{}ptr}
handle that closes exactly once and skips writing if \code{fopen} fails; a
\code{printf} with three \code{\%le} and one argument and an unsynchronised
counter in its OpenMP loop were also fixed. Every other \code{fopen} in the
code now checks for failure.
\end{response}
\status{Fixed (\commit{ac039a2}, \commit{538ed73}).}

\begin{reviewer}{Email 10/15; Summary}
Possible \code{log()} of a non-positive value in the Box--Cox transform
(\code{ApplyBoxCoxTransform}).
\end{reviewer}
\begin{response}
Confirmed for zero values (negative values were already skipped; zeros gave
$-\infty$). Any constituent with a zero or negative value is now left
untransformed, and the analysis lists the constituents it skipped.
$\lambda = 0$ is handled as the log case; for $|\lambda| < 0.01$ the
transform is computed as $\mathrm{expm1}(\lambda \log x)/\lambda$ so it keeps
full precision; elsewhere the formula is unchanged, so optimal $\lambda$
values do not move. Two related faults were fixed: outlier detection used
its own copy of the formula, which scored skipped constituents from $x-1$
and so shifted their z-scores by $1/\sigma$; and the (unused) Box--Cox
log-likelihood took the log of a sum that was always negative.
\end{response}
\status{Fixed (\commit{ac039a2}). Tests: \code{boxcox\_\allowbreak{}positiveData},
\code{boxcox\_\allowbreak{}zeroAndNegativeData}, \code{boxcox\_\allowbreak{}lambdaZeroAndNearZero},
\code{boxcox\_\allowbreak{}logLikelihoodIsFinite}.}

\begin{reviewer}{Email 10/15; Summary}
Widespread uninitialised variables flagged by clang-tidy and cppcheck.
\end{reviewer}
\begin{response}
All 128 findings of clang-tidy's \code{cppcoreguidelines-\allowbreak{}init-variables}
and \code{cppcoreguidelines-\allowbreak{}pro-type-member-init} were fixed: in-class
initialisers for members, and NaN for floating-point locals so that any
read-before-write is visible rather than silent. One was a real defect: the
stepwise DFA could report an unassigned Wilks' $\lambda$ when no element
improved the fit (it now reports NaN). These two checks are errors in the CI
clang-tidy job, so new instances fail the pipeline.
\end{response}
\status{Fixed (\commit{ac039a2}); enforced in CI.}

% ====================================================================
\group{Reproducibility and random numbers}

\begin{reviewer}{Email 10/15, item 1; Summary; Tests, Test 1}
The MCMC is non-reproducible: identical inputs give materially different
outputs (Bank upper credible bound 0.47--0.67 over three runs), there is no
seed control and no recorded seed. ``Consider adding a way to provide a
user-defined seed to allow for testing.''
\end{reviewer}
\begin{response}
Done. CMB Bayesian, CMB Bayesian-Batch and Test CMB Bayesian have a new
\emph{Random seed} field (blank = drawn from the operating system). The seed
used is shown in the name of the results (``MCMC results for `X' (random
seed $N$)''), stored with the project and in the command-line report, and
written on the first line of the samples file. The same seed and settings
give byte-identical results. The generator is \code{std::\allowbreak{}mt19937\_\allowbreak{}64} seeded
through \code{std::\allowbreak{}seed\_\allowbreak{}seq}; the uniform and normal values are computed
from its raw output in \code{include/\allowbreak{}randomstream.\allowbreak{}h} rather than with
\code{std::\allowbreak{}*\_\allowbreak{}distribution}, whose algorithms differ between standard
libraries, so a recorded seed reproduces a run on Linux, macOS and Windows.
\end{response}
\status{Fixed (\commit{cdfa4ea}). Tests: \code{mcmc\_\allowbreak{}sameSeedSameResult},
\code{mcmc\_\allowbreak{}differentSeedDifferentResult}, \code{random\_\allowbreak{}streamsAreReproducible}, \code{random\_\allowbreak{}seedParsing}.}

\begin{reviewer}{Email 10/15, item 2; Summary; RABBIT\_HOLE\_rand\_and\_threads.md}
\code{rand()}/\code{srand()} called inside an OpenMP parallel region is a data
race that also degrades the generator's statistical properties.
\end{reviewer}
\begin{response}
Kiro was right that this was a race. Each chain now has its own stream,
\code{RandomStream(seed, chain)}, used for its random start and every step,
so no generator is shared between threads and the result does not depend
on scheduling: the samples are identical with 1, 3 or 8 threads. I audited
every OpenMP region: the MCMC step loop also wrote a shared cached maximum
from every thread (now a named critical section), and the GA loop's
counter is now atomic. Every other \code{rand()}, \code{srand()} and
clock-seeded generator in the code now draws from a per-thread stream with
a fixed seed, and Error Analysis (which removes random samples) now has its
own seed field like the MCMC. clang-tidy's \code{concurrency-mt-unsafe} check
is now an error in CI, so \code{rand()} cannot creep back in.

On the details of \code{RABBIT\_\allowbreak{}HOLE\_\allowbreak{}rand\_\allowbreak{}and\_\allowbreak{}threads.\allowbreak{}md} (your instinct
not to bet the house was sound): the conclusion and the remedy are right,
but the mechanism is overstated for the platforms SedSAT3 runs on. The
\code{srand} block only ran on Windows (\code{\#ifdef WIN64}), where the
MSVC runtime keeps \code{rand()} state per thread, and glibc's \code{rand()}
takes a lock, so ``lost updates'' were theoretical. The real problems were
that the order of draws depended on thread scheduling, Windows reseeded
from the clock every iteration, MSVC's \code{RAND\_\allowbreak{}MAX} of 32767 gave
uniforms only 15 bits of resolution, and Linux and Windows produced
different sequences. Its suggested fix also needed two corrections, both
made: \code{std::\allowbreak{}normal\_\allowbreak{}distribution} is not portable across standard
libraries, and streams must be per \emph{chain}, not per thread, for results
to be independent of the thread count.
\end{response}
\status{Fixed (\commit{cdfa4ea}). Test: \code{mcmc\_\allowbreak{}independentOfThreadCount}.}

\begin{reviewer}{Tests, Test 1 (convergence) and your note on it}
Increasing samples from 500 to 3000 reduced, but did not remove, the
run-to-run spread, and the software surfaces no convergence diagnostics
(R-hat, effective sample size, trace plots). [Your note: ``that conclusion
is a little strong.'']
\end{reviewer}
\begin{response}
Agreed on both counts. Part of the spread came from a units issue in the
dialog: \emph{Number of samples} and the burn-in count samples across
\emph{all} chains, so 3000 samples with 8 chains is only 375 steps per chain;
at 40\,000 samples with 20\,000 burn-in the validation targets in our test
project are recovered at 85--99\% to the correct source, in agreement with
Levenberg--Marquardt. Now that runs are reproducible, sampling noise and
non-convergence can at least be separated by re-running with the same seed
and more samples. Exposing R-hat, effective sample size and trace plots is a
feature I plan for a later version; whether the current defaults are
adequate is for the domain review.
\end{response}
\status{Partly addressed (reproducibility); diagnostics deferred; domain question handed off.}

\begin{reviewer}{Tests, LM contrast}
The deterministic Levenberg--Marquardt method reproduces exactly, which
shows the variation is specific to the stochastic path.
\end{reviewer}
\begin{response}
On checking, LM was less deterministic than it looked: its random starting
point came from unseeded \code{rand()}, so it repeated only while nothing
reseeded the generator (the MCMC did, on Windows) and differed between Linux
and Windows. It now draws its starting point from a fixed-seed stream created
for each solve, so it is repeatable everywhere. Most results changed by
less than $10^{-5}$, but for one target sample in the domain-review project
they changed by up to 0.06, showing that the LM solution for that sample
depends on its starting point. I have flagged this for the domain review.
\end{response}
\status{Fixed (\commit{cdfa4ea}); regression baseline updated with explanation.}

% ====================================================================
\group{Input validation}

\begin{reviewer}{Email 10/15, item 3; Tests, Test 4b; Email 10/16}
A bare ``NaN'' passes \code{QString::\allowbreak{}toDouble} and is stored raw; whether it
then affects a result depends on whether each consumer calls
\code{removeNaNs()} or checks \code{isfinite}, which is inconsistent.
``Have Claude help create some unit tests that would go after the odd NaN /
legitimately negative / otherwise illegal values.''
\end{reviewer}
\begin{response}
There is now one parser, \code{include/\allowbreak{}numericparsing.\allowbreak{}h}
(\code{parseNumericCell} and overloads for spreadsheet and project-file
values), used by the Excel import, project loading, the element table and
the command line. It refuses empty cells, NaN / Inf / Infinity in any case,
overflow, text, and decimal commas or thousands separators, each with a
specific reason. A refused value stops the import with sheet, row, column
and value (Excel) or group, sample and constituent (project file) — the same
treatment empty and non-numeric cells already had, since SedSAT3 has no
missing-data value. Two findings from tracing the behaviour you observed:
the ``Ignored NaN, Inf, or -Inf value'' message came from Qt Charts while
plotting, so the NaN \emph{had} been stored; and project loading turned a
missing value or a NaN (which Qt saves as \code{null}) into 0, and silently
dropped a sample that failed to load. Both now stop the load with the
location. The only ad-hoc data guard downstream (the Box--Cox optimiser's
\code{min == min} test) uses a shared helper; the other \code{isfinite}
uses guard computed quantities, not data, and \code{removeNaNs()} has no
callers. The import is now a separate, dialog-free function
(\code{excelimport.cpp}), so it is tested on real workbooks; it also refuses
an empty header cell, a row without a sample name and a repeated sample name,
each of which used to drop data silently. The unit tests you suggested cover 8 valid formats and 14 illegal
inputs, spreadsheet and JSON values, a saved-and-reloaded NaN, zero and
negative data in Box--Cox, and negative values in the source groups.
\end{response}
\status{Fixed (\commit{68be7ca}). Tests: \code{parser\_\allowbreak{}acceptsNumbers}, \code{parser\_\allowbreak{}rejects},
\code{parser\_\allowbreak{}variantAndJson}, \code{boxcox\_\allowbreak{}zeroAndNegativeData}, \code{negativeCheck\_\allowbreak{}*}.}

\begin{reviewer}{Tests, Test 5}
Project persistence round-trips correctly (data, tracer exclusions and MCMC
results survive save and reload).
\end{reviewer}
\begin{response}
Thank you for checking. Saved projects remain backward compatible: no
serialised key was renamed, and projects with only finite values load
exactly as before.
\end{response}
\status{No change needed.}

% ====================================================================
\group{Tests, CI, the genetic algorithm, and hygiene}

\begin{reviewer}{Email 10/15, item 4; Summary}
There appear to be no automated tests or CI beyond a docs workflow; add at
least one end-to-end regression against a known result and run the
static-analysis tasks in CI.
\end{reviewer}
\begin{response}
Added. \code{tests/\allowbreak{}unit/\allowbreak{}} has 49 Qt Test checks covering every defect above.
\code{tests/\allowbreak{}regression/\allowbreak{}} runs the deterministic analyses (Box--Cox, outliers,
Kolmogorov--Smirnov, distribution fitting, ANOVA, correlation and covariance,
bracketing, DFA, EDP, Levenberg--Marquardt batch) on two example projects
and compares them with a stored baseline at a relative tolerance of
$10^{-9}$; the output of the code before the fixes is kept alongside, so every
change the fixes made is traceable. \code{.\allowbreak{}gitlab-ci.\allowbreak{}yml} builds with CMake
(GUI, command line, tests), with the GA enabled, and with the qmake
projects; runs the unit and regression tests; runs clang-tidy (initialisation
checks as errors) and cppcheck; has a manual Windows job for a Windows
runner; and publishes the Doxygen documentation to GitLab Pages. The
existing GitHub documentation workflow is kept for the development mirror.
\end{response}
\status{Fixed (\commit{2f8c463}).}

\begin{reviewer}{Email 10/15, item 5; Summary; Tests, Test 6}
The genetic algorithm is advertised in the README and code.json but is in
no menu and does not compile (missing \code{StringOP.\allowbreak{}h}); either wire it up
or drop the claim. \code{MCMC.\allowbreak{}h} includes \code{GA.\allowbreak{}h}, so the GA header is
not cleanly deletable.
\end{reviewer}
\begin{response}
The GA works and is deliberately not offered: Levenberg--Marquardt solves
the same problem more efficiently, and two deterministic methods would
confuse users. It is now kept in the source but excluded from the default
build (\code{CONFIG+=ga} with qmake, \code{-DENABLE\_\allowbreak{}GA=ON} with CMake); both
settings build and CI builds both. \code{MCMC.\allowbreak{}h} no longer includes
\code{GA.\allowbreak{}h}. The README and code.json no longer list it as a method, the
README says it is included but disabled and why, and the command-line program
and its manual leave out its commands. The missing \code{StringOP.\allowbreak{}h} was only
included by the stale \code{GA.\allowbreak{}cpp} (Comment 2.7).
\end{response}
\status{Fixed (\commit{2f8c463}). CI job \code{build:linux-ga}.}

\begin{reviewer}{Email 10/15, item 6; Summary}
A dead submodule URL blocks fresh clones; stray build artifacts; spelling in
identifiers (\code{purturbation}, \code{FileBrowserPushButtom}, ``samle'').
\end{reviewer}
\begin{response}
The QXlsx submodule URL works (it redirects to the upstream
\code{QtExcel/\allowbreak{}QXlsx}, and the pinned commit is fetchable); it now points at
the upstream directly. The failing clone was most likely the
\code{Utilities} submodule, which is gone. A fresh
\code{git clone --recursive} of the branch builds. Removed: \code{wasm\_\allowbreak{}bin/\allowbreak{}}
build output, the stale \code{CMBSource.\allowbreak{}pro} (which listed files that do not
exist), a stray \code{gitignore} copy and an already-applied script; the
Linux build no longer links a prebuilt \code{libQXlsx.\allowbreak{}a} but compiles QXlsx
from the submodule; \code{.\allowbreak{}gitignore} covers build folders, IDE files and
binaries. Spelling: \code{FileBrowserPushButton}, \code{perturbation} (the
old MCMC setting names are still accepted), ``sample'', ``Description''
(the misspelled form appeared as a label in the GUI) and
\code{sizeconstituent}.
\end{response}
\status{Fixed (\commit{2f8c463}).}

\begin{reviewer}{Summary, optional / polish}
Three coexisting build systems with no documented source of truth; repo
layout mid-transition (non-GUI core at the repository root).
\end{reviewer}
\begin{response}
The qmake projects now share one file, \code{sedsat3\_\allowbreak{}core.\allowbreak{}pri}, for the
analysis sources, libraries and options; CMake builds the GUI, the command
line program and the tests from the same source list; the Visual Studio
project is kept for Windows developers; the README documents all three. I
have left the move of the root sources under \code{src/\allowbreak{}} for a separate
change, since it is large and mechanical and would make this review harder
to follow.
\end{response}
\status{Improved (\commit{2f8c463}); layout change deferred.}

\begin{reviewer}{Email 10/15; Email 10/16; Linter}
Compile with multiple compilers with warnings ``turned up to 11'' and
systematically reduce them; the more significant warnings generally point at
minor defects.
\end{reviewer}
\begin{response}
I did this with GCC~13 (and clang-tidy from LLVM~20). Unique first-party
warnings in the GUI build went from 244 to 100 with \code{-Wall -Wextra}, and
from 618 to 448 with \code{-Wall -Wextra -Wpedantic -Wshadow -Wconversion}.
As you predicted, the significant ones were real defects: two operators that
called themselves (stack overflow), a constructor that wrote past the end of
an empty matrix, a function returning a reference to a local, printf format
mismatches, invalid escapes, and the missing returns. Every
\code{-Wreturn-type}, \code{-Winfinite-recursion}, \code{-Wreturn-local-addr},
\code{-Wformat}, \code{-Wmaybe-uninitialized}, \code{-Wmisleading-indentation},
\code{-Wreorder}, unused-variable/parameter and copy-constructor warning is
gone. The remainder are implicit integer/floating conversions
(\code{-Wconversion} 293, \code{-Wsign-compare} 100, \code{-Wfloat-conversion}
20) and \code{-Wshadow} 35 (mostly constructor parameters named after
members); I prefer to change the numerical conversions one at a time with
the regression tests rather than in bulk.
The clang-tidy findings in \code{linter\_\allowbreak{}output.\allowbreak{}txt} are answered in
Section~9.
\end{response}
\status{Largely fixed (\commit{2f8c463}).}

% ====================================================================
\group{Further defects found and fixed during the response}

\begin{reviewer}{Found while addressing the review}
Defects not in the review, found while fixing the items above.
\end{reviewer}
\begin{response}
\begin{itemize}[nosep,leftmargin=*]
\item ANOVA read the ``Use only selected samples/elements'' options and then
  ignored them (it always used the selected samples and every element).
\item \code{NegativeValueCheck}, \code{GrandMean} and \code{LumpAllProfileSets}
  skipped the source sorting after the target (Comment 2.4); ANOVA's total sum
  of squares left that source out.
\item \code{CMBTimeSeries::\allowbreak{}operator=} and \code{CVector\_\allowbreak{}arma::\allowbreak{}operator*=}
  called themselves (stack overflow on use).
\item \code{CMatrix\_\allowbreak{}arma(const CVector\_\allowbreak{}arma\&)} built a temporary instead of
  sizing the matrix, then wrote past its end.
\item \code{CMatrix\_\allowbreak{}arma\_\allowbreak{}sp::\allowbreak{}get} returned a reference to a local, so writes
  to sparse matrices were lost.
\item \code{RangeSet}'s copy constructor copied nothing; the copy constructors
  of several result types dropped their notes and axis labels.
\item The MCMC's shared prediction rows were written from every thread
  (Comment 3.2).
\item \code{TimeSeries::\allowbreak{}addPoint} did not compare the first two intervals,
  so an uneven series could be treated as evenly spaced.
\item \code{gamma()} clashed with \code{::\allowbreak{}gamma} under Clang.
\end{itemize}
\end{response}
\status{Fixed (\commit{ac039a2}, \commit{2f8c463}).}

\begin{reviewer}{Effect on results}
How do the fixes change the results the software produces?
\end{reviewer}
\begin{response}
The regression suite was run on the code before the fixes and after; every
difference is explained in \code{docs/\allowbreak{}review/\allowbreak{}REVIEW\_\allowbreak{}RESPONSE.\allowbreak{}md}. In short:
distribution fitting now succeeds where it crashed; outlier z-scores of
constituents with negative values (C13; N15 in Forest) shift by a constant
$1/\sigma$; Box--Cox and outlier analyses list the constituents they skip;
Levenberg--Marquardt moves by less than $10^{-5}$ for most samples and up to
0.06 for one (Comment 3.4); and ANOVA p-values change, some across 0.05,
because the selection options now take effect and the previously dropped
source group is included. Box--Cox $\lambda$, Kolmogorov--Smirnov,
correlation and covariance, bracketing, DFA and EDP are unchanged. The ANOVA
and Levenberg--Marquardt changes have been passed to the domain reviewers.
\end{response}

% ====================================================================
\group{Windows build and the portable package}

\begin{reviewer}{Email 10/16; BUILDING\_ON\_WINDOWS\_WITH\_PIXI.md}
The reviewed copy had no viable Windows compilation path; Kiro created one
with pixi (conda-forge libraries driven from \code{pixi.\allowbreak{}toml}, with
\code{pixi run lint} for the strict-warning build). ``This may or may not be
useful.''
\end{reviewer}
\begin{response}
It is useful, and since it is how you build, I want it to work out of the box
at the re-review. Several of the pitfalls your guide works around are now gone
from the repository: the QXlsx URL is fixed, \code{Utilities} is no longer an
(empty) submodule, the hard-coded \code{\#define ARMA\_\allowbreak{}USE\_\allowbreak{}SUPERLU} that broke
the win-64 link has been removed (SuperLU is now used only if the installed
Armadillo provides it), MSVC gets \code{/\allowbreak{}utf-8} from CMake (your C4828), and
the duplicate \code{USE\_\allowbreak{}QCHARTS} definition is guarded (your C4005). The guide
itself needs three adjustments to match the repository: its \code{lint} task
lists the removed \code{src/\allowbreak{}GA/\allowbreak{}GA.\allowbreak{}cpp}; its \code{.\allowbreak{}clang-tidy} would replace
the repository's stricter one (which CI uses, with the initialisation and
\code{rand()} checks as errors); and \code{deploy-win} uses
\code{--compiler-runtime} and LLVM's OpenMP DLLs, whereas MSVC builds need
\code{vcomp140.\allowbreak{}dll} (see \code{docs/\allowbreak{}PACKAGE-Windows-Portable.\allowbreak{}md}).
A corrected \code{pixi.toml} is now in the repository. It drives the
existing CMake build and has tasks \code{build}, \code{test} (unit and
regression tests), \code{lint} (with the repository's \code{.clang-tidy}),
\code{cppcheck} (with your \code{std.cfg} workaround), \code{warnings} (an
MSVC \code{/W4} build), \code{deploy-win} and \code{run-app}.
\mine{Run on Windows: pixi run build, test, lint, cppcheck and warnings; then
replace this note with the results, including the MSVC /W4 and cppcheck counts} The Visual Studio project now compiles
the vendored QCustomPlot instead of a copy outside the repository; it still
links a prebuilt QXlsx library, which I will change to the submodule as the
qmake and CMake builds already do. The portable, no-administrator Windows
package (unzip and run) will be built from the merged code and sent to the
reviewers.
\end{response}
\status{Build fixes done; \code{pixi.\allowbreak{}toml} pending a Windows test.}

% ====================================================================
\group{Detailed observations (CODE\_REVIEW\_OBSERVATIONS.md)}

\begin{reviewer}{Observations, \S1--16}
The full review log, ``nuggets of gold mixed with Kiro fluff''.
\end{reviewer}
\begin{response}
I checked every concrete finding against the code. Most were already
addressed by the changes above; the rest are listed here with what was
done. A few did not hold up when checked, and I note them so the record is
accurate, not to diminish the review: they were all worth looking at.
\end{response}

{\small
\begin{longtable}{@{}p{0.09\linewidth}p{0.38\linewidth}p{0.46\linewidth}@{}}
\toprule
\S & Finding & Response\\
\midrule
\endhead
1 & No command-line entry point & A command-line program, \code{sedsat3-cli}, runs any analysis from a JSON script; the regression tests use it.\\
1, 9 & GA listed as a method; GA not built & Fixed (Comment 5.2).\\
2 & Three build systems, no source of truth & Improved (Comment 5.4).\\
2 & \code{CMBSource.\allowbreak{}pro} hard-codes a personal path & File removed.\\
2 & The same kind of paths in \code{SedSat3.\allowbreak{}vcxproj} & QCustomPlot now from \code{thirdparty/\allowbreak{}}; the prebuilt QXlsx still to change (Section 7).\\
2 & Hygiene script not run; stray \code{gitignore} & Both removed.\\
2 & \code{CMB\_\allowbreak{}Source.\allowbreak{}aip} still tracked & Kept on purpose (the installer) and rebranded to the USGS.\\
3a & QXlsx submodule URL dead & Fixed (Comment 5.3).\\
3b & \code{Utilities} is a personal-namespace submodule & Brought into the repository (Comment 1.2).\\
3b & Hard-coded \code{ARMA\_\allowbreak{}USE\_\allowbreak{}SUPERLU} breaks Windows & Fixed (Section 7).\\
3b & No Windows build path in the repository & MinGW guide, Visual Studio project, CMake, a manual Windows CI job, and \code{pixi.toml} (Section 7).\\
3c & qcustomplot6 version not documented & \code{THIRD\_\allowbreak{}PARTY\_\allowbreak{}LICENSES.\allowbreak{}md} (2.1.1).\\
4, 4a & README truncated; advertises sedsat.org, GA, GitHub & Rewritten (Comments 1.5, 1.6).\\
4 & Dependency versions unpinned & README and \code{THIRD\_\allowbreak{}PARTY\_\allowbreak{}LICENSES.\allowbreak{}md} give versions (Qt 6.2+, 6.8 tested, \ldots).\\
4a & ``CMB''-era artifacts & Removed (only icon file names remain).\\
5 & No automated tests or CI & Fixed (Comment 5.1).\\
5, 9 & \code{testmcmc} is a non-functional stub & It is the model behind ``Test CMB Bayesian'', and is now the fixture of the MCMC reproducibility tests.\\
6 & MSVC C4828 (non-UTF-8 bytes) & Four files were Latin-1; all sources are now UTF-8, and CMake passes \code{/\allowbreak{}utf-8} to MSVC.\\
6 & MSVC C4005, C4834 & C4005: the \code{USE\_\allowbreak{}QCHARTS} and \code{NOMINMAX} redefinitions are now guarded. C4834 not reproduced (no MSVC here); will check with the first MSVC \code{/\allowbreak{}W4} run.\\
9 & Non-GUI core at the root; mixed header layout & Deferred to a separate change (Comment 5.4).\\
9 & \code{FileBrowserPushButtom} misspelling & Fixed.\\
9 & Duplicate \code{FilePushButton} class & Unused; removed.\\
9 & Overlapping chart layers (Qt Charts and QCustomPlot) & Acknowledged; to be resolved together with the GPL-dependency decision (Comment 1.3).\\
9 & Orphan \code{drand.\allowbreak{}cpp}; empty \code{GA\_\allowbreak{}impl} files & Removed.\\
10 & Clock-seeded \code{srand}; no seed control or record & Fixed (Comment 3.1).\\
10 & Every thread's default stream starts from the same seed & Documented: parallel code must use its own seeded stream, as the MCMC does; no parallel code uses the default stream.\\
11.1--11.2 & Platform-specific seeding; \code{rand()} race in OpenMP & Fixed (Comment 3.2).\\
11.3 & GA file I/O inside an OpenMP critical section & Correct but slow; left as is because the GA is off by default.\\
11.4 & \code{qDebug} in the MCMC loop & Removed.\\
11 & OpenMP use not documented & The chains run in parallel and results do not depend on the thread count; noted in \code{include/\allowbreak{}randomstream.\allowbreak{}h} and the README.\\
12 & Project declares clang-tidy/cppcheck in \code{pixi.\allowbreak{}toml} & There was no \code{pixi.\allowbreak{}toml} in the repository; this was your local environment. Both tools now run in CI.\\
12a.1 & Null dereference in \code{CBinary::\allowbreak{}operator[]} & Fixed (Comment 2.6).\\
12a.2 & Writes to a closed \code{FILE*} & Fixed (Comment 2.7).\\
12a.3 & Missing return in both \code{CheckNegativeElements} overloads & Fixed; in the reviewed snapshot only the map overload lacked a return, and its result was not used, so the guard itself worked (Comment 2.4).\\
12a.4 & \code{log()} of a non-positive value in \code{ApplyBoxCoxTransform} & Fixed; the always-negative \code{log} argument was in \code{CalculateBoxCoxLogLikelihood} (Comment 2.8).\\
12b & \code{Interface} limits uninitialised; copy operations incomplete & Fixed: members initialised, and the copy constructor and assignment now copy every member, so copies keep their notes (outlier and bracketing flags) and highlight limits.\\
12b & Uninitialised members of \code{Observation}, \code{Parameter}, \code{element\_\allowbreak{}information} & Fixed (Comment 2.9).\\
12b, 12d & Performance and implicit-widening notes & See Section 9.\\
12d & Missing \code{StringOP.\allowbreak{}h}; GA headers use incomplete types & Fixed; GA.hpp now includes what it uses, and the GA builds in CI.\\
12d & \code{MCMC.\allowbreak{}h} includes \code{GA.\allowbreak{}h} & Fixed.\\
13 & GPL-3.0 outbound license; GPL dependencies & Comment 1.3.\\
13 & GSL used for a single \code{gsl\_\allowbreak{}multifit\_\allowbreak{}wlinear} call & GSL is used more widely (CDFs, statistics, matrices in about ten files), so replacing it is not a one-call change; another reason to prefer the inventory approach.\\
13, 13a & No per-file headers; AGPL headers in \code{Utilities} & Fixed (Comments 1.2, 1.7).\\
14, 14a & Metadata and installer point outside the USGS & Fixed (Comment 1.6); the parallel sedsat.org distribution is for the release coordinator.\\
14 & \code{createdmg.\allowbreak{}sh} version stale & Updated to 1.1.7.\\
16 & Spelling; add codespell to CI & Spelling fixed (Comment 5.3); codespell now runs in CI (reported, not yet enforced).\\
8.4 & Distribution Fitting crash caused by a missing overload & Fixed; the cause was an implicit conversion (Comment 2.1).\\
8.6 & ``Include all elements'' inverted & Fixed (Comment 2.2).\\
\bottomrule
\end{longtable}
}

\begin{reviewer}{Found while checking the observations}
Further defects found while verifying the findings above.
\end{reviewer}
\begin{response}
\begin{itemize}[nosep,leftmargin=*]
\item \code{ObservedDataforSelectedSample\_\allowbreak{}Isotope} read the base element from
  an empty sample name in one branch, dereferencing a null pointer.
\item The two-sample t-test computed its degrees of freedom as an unsigned
  value that wrapped around with fewer than two samples in a group; it now
  reports NaN.
\item The distribution curve for uniform and Dirichlet priors consisted of
  NaN points, and a one-point curve divided by zero.
\item The results table compared the \emph{time} rather than the value with
  the upper highlight limit.
\item An MCMC initialisation routine computed a step from $-\mathrm{low} - \mathrm{high}$
  instead of $\mathrm{high} - \mathrm{low}$.
\item \code{CGA::\allowbreak{}operator=} took a non-const reference and returned by value.
\item The Excel import silently dropped data after an empty header cell or a
  row without a sample name, and dropped repeated sample names (found while
  testing your Test 4a).
\end{itemize}
All are fixed. Unit tests were added for the uniform curve, for copies
keeping their notes, and for the Excel import.
\end{response}

% ====================================================================
\group{clang-tidy findings (linter\_output.txt)}

\begin{reviewer}{Linter}
\code{pixi run lint} output: clang-tidy over 13 translation units with
\code{bugprone-*}, \code{clang-analyzer-*}, \code{concurrency-*},
\code{cppcoreguidelines-\allowbreak{}init-variables}, \code{-pro-type-member-init},
\code{performance-*} and \code{modernize-*} checks (101 first-party
diagnostics, 8 of them errors).
\end{reviewer}
\begin{response}
Every error and every finding that indicated a defect is fixed or no longer
applies:
\end{response}

{\small
\begin{longtable}{@{}p{0.42\linewidth}p{0.08\linewidth}p{0.44\linewidth}@{}}
\toprule
Check & Count & Status\\
\midrule
\endhead
\code{cppcoreguidelines-\allowbreak{}init-variables} & 35 & Fixed (Comment 2.9); one in the old \code{GA.\allowbreak{}cpp} was a false positive.\\
\code{cppcoreguidelines-\allowbreak{}pro-type-member-init} & 13 & Fixed with in-class initialisers.\\
\code{clang-diagnostic-error}: \code{StringOP.\allowbreak{}h} not found & 1 & File removed (Comment 2.7).\\
\code{clang-diagnostic-error}: incomplete types in \code{GA.\allowbreak{}hpp} & 7 & Fixed: \code{GA.\allowbreak{}hpp} includes \code{progressreporter.\allowbreak{}h} itself, and its \code{Q\_\allowbreak{}version} guards (a macro defined nowhere) are now \code{Q\_\allowbreak{}GUI\_\allowbreak{}SUPPORT}.\\
\code{clang-analyzer-cplusplus.\allowbreak{}NewDeleteLeaks} (progress windows) & 5 & Fixed: progress windows are owned by a scoped object.\\
\code{concurrency-mt-unsafe} (\code{rand}) & 2 & Fixed (Comment 3.2); now an error in CI.\\
\code{clang-analyzer-deadcode.\allowbreak{}DeadStores} (ANOVA options) & 2 & Fixed: the options were ignored (Section 6).\\
\code{UndefinedBinaryOperator}, \code{CallAndMessage} & 4 & Fixed (initialisation, uniform curve).\\
\code{UndefReturn}, \code{modernize-use-nullptr} (\code{CBinary}) & 2 & Fixed (Comment 2.6).\\
\code{bugprone-unhandled-self-assignment} & 1 & Fixed (\code{CGA::\allowbreak{}operator=}).\\
\code{bugprone-branch-clone} (GA parameter mapping) & 1 & Fixed: \code{apply\_\allowbreak{}to\_\allowbreak{}all} is false for every parameter in SedSAT3 (a single model), a remnant of a multi-series design; the identical branches were collapsed.\\
\code{bugprone-integer-division} & 1 & Intended (the result is bounded below by 1).\\
\code{bugprone-float-loop-counter}, \code{FloatLoopCounter} & 6 & Grid loops over $x$ and $\lambda$; a zero step is now impossible (Comment 2.1). Kept.\\
\code{bugprone-throwing-static-initialization} & 1 & False positive (\code{atan(1.0)}).\\
\code{bugprone-\allowbreak{}implicit-\allowbreak{}widening-\allowbreak{}of-\allowbreak{}multiplication-\allowbreak{}result} & 4 & Fixed with explicit casts (11 across the whole code).\\
\code{performance-\allowbreak{}unnecessary-value-param}, \code{performance-for-range-copy} & 15 & Fixed across the whole code (83 parameters now passed by \code{const\&}, 11 loops no longer copy); both checks are now enabled in CI.\\
\bottomrule
\end{longtable}
}

Run over all first-party code with the repository's configuration (your
check set, plus the performance checks), clang-tidy now reports 13 findings:
7 grid loops over a floating-point counter (intended), 3 destructors that
call Armadillo's \code{clear()}, 2 false positives on Qt-parented labels, and
one readable \code{if}/\code{else} chain.

% ====================================================================
\group{Reviewer GUI test plan (REVIEWER\_TEST\_PLAN.md)}

\begin{reviewer}{REVIEWER\_TEST\_PLAN.md and TEST\_RESULTS.md}
Tests 1--7 and the GUI observations G1--G3.
\end{reviewer}

{\small
\begin{longtable}{@{}p{0.2\linewidth}p{0.74\linewidth}@{}}
\toprule
Test & Status and automated test\\
\midrule
\endhead
1. MCMC reproducibility & Fixed; \code{mcmc\_\allowbreak{}sameSeedSameResult}, \code{mcmc\_\allowbreak{}differentSeedDifferentResult}. Convergence diagnostics deferred (Comment 3.3).\\
1b. LM determinism & Made genuinely repeatable (Comment 3.4); regression suite.\\
2. Sum to one, non-negative & Now tested for both parameterisations: \code{levenbergMarquardt\_\allowbreak{}contributionsSumToOne}. The Bayesian contributions satisfy it by construction (the last is one minus the others).\\
3. LM vs Bayesian & Not run by the review; left to the domain review.\\
4a. Required column removed & Now tested (\code{excelImport}): a column removed from one sheet is refused, naming the two sheets; removed from every sheet, the workbook imports without it. Testing this found three ways an import silently dropped data — an empty header cell (every column after it), a row without a sample name (every sample below it), and a repeated sample name (the second sample) — all now refused with their location.\\
4b. NaN / non-numeric cell & Fixed (Comment 4.1); parser tests and \code{excelImport}.\\
4c. Negative values & Fixed (Comment 2.4); \code{negativeCheck\_\allowbreak{}*}.\\
4c addendum & Message improved (Comment 2.5).\\
4d. Single-sample source & Now tested: \code{levenbergMarquardt\_\allowbreak{}singleSampleSource}, \code{distributionFitting\_\allowbreak{}conductor}.\\
5. Round trip & Passed in your test; unchanged.\\
6. GA advertised & Fixed (Comment 5.2).\\
7. RNG thread safety & Fixed; \code{mcmc\_\allowbreak{}independentOfThreadCount}.\\
Test 7 (results). Distribution Fitting crash & Fixed (Comment 2.1).\\
G1--G3 & Fixed (Comments 2.2, 2.3).\\
\bottomrule
\end{longtable}
}

The test plan's notes that ``SedSAT3 has no CLI'' and that it is launched with
\code{pixi run run-app} are now out of date; Tests 1 and 7 can be scripted with
\code{sedsat3-cli}.

% ====================================================================
\group{USGS software release checklist}

\begin{reviewer}{software-release-checklist.pdf}
The USGS checklist for open-source and scientific-information-product
releases.
\end{reviewer}
\begin{response}
The repository items are in place; the remaining items are decisions or
process steps for the release coordinator, Cara and me:
\end{response}

{\small
\begin{longtable}{@{}p{0.36\linewidth}p{0.2\linewidth}p{0.38\linewidth}@{}}
\toprule
Item & Status & Remaining\\
\midrule
\endhead
README.md, DISCLAIMER.md (provisional) & Done & Placeholders: repository URL, manual, citation details.\\
LICENSE.md & Placeholder & License choice (Comment 1.3).\\
code.json & Done, placeholders & Version, contact, DOI, URLs; \code{name} to match the GitLab project; \code{downloadURL} to point at \code{main}; \todo{ask whether the AI tag applies to AI-assisted development}.\\
Project in a USGS group (not a personal namespace) & Not done & Release coordinator: move from \code{cphipps/\allowbreak{}}.\\
Default branch is what is released & In progress & Merge request !4 into \code{main}.\\
Administrative security review of every commit & Process & Security reviewer; the history includes old binaries and the AGPL-era files; decide whether to publish or rewrite it.\\
IPDS record; reviews logged & Process & Cara: log this review, this response and the re-review.\\
DOI reserved & Process & Release coordinator.\\
Release-candidate branch named for the version & Not done & After the version is chosen (semantic version, no ``v'').\\
Domain review & Open & LM start sensitivity, MCMC convergence, OM/size correction producing negatives, isotopic tracers, ANOVA changes.\\
Final changes on the release branch (approved disclaimer, tagged URLs) & Not done & At release.\\
\bottomrule
\end{longtable}
}

\end{document}
